\documentclass[11pt]{article}
\usepackage[a4paper,margin=28mm]{geometry}
\usepackage[T1]{fontenc}
\usepackage[utf8]{inputenc}
\usepackage{lmodern}
\usepackage{amsmath,amssymb,amsthm,mathtools}
\usepackage{booktabs,longtable,array}
\usepackage{microtype}
\usepackage[colorlinks=true,linkcolor=blue,citecolor=blue,urlcolor=blue]{hyperref}
\hypersetup{pdfauthor={Anonymous},
pdftitle={Sharp rank bounds and all-parameter tree-network separation},
pdfsubject={Unrefereed candidate manuscript}}
\usepackage{enumitem}
\setlist{nosep}
\newtheorem{theorem}{Theorem}[section]
\newtheorem{lemma}[theorem]{Lemma}
\newtheorem{proposition}[theorem]{Proposition}
\newtheorem{corollary}[theorem]{Corollary}
\theoremstyle{definition}
\newtheorem{definition}[theorem]{Definition}
\theoremstyle{remark}
\newtheorem{remark}[theorem]{Remark}
\DeclareMathOperator{\rank}{rank}
\DeclareMathOperator{\diag}{diag}
\DeclareMathOperator{\tr}{tr}
\DeclareMathOperator{\dist}{dist}
\DeclareMathOperator{\TV}{TV}
\newcommand{\one}{\mathbf 1}
\newcommand{\R}{\mathbb R}
\newcommand{\E}{\mathbb E}
\newcommand{\Prob}{\mathbb P}
\newcommand{\Flat}{\operatorname{Flat}}
\newcommand{\ReleaseDOI}{\href{https://doi.org/10.5281/zenodo.23140370}{10.5281/zenodo.23140370}}
\newcommand{\doi}[1]{\href{https://doi.org/#1}{doi:\nolinkurl{#1}}}
\title{Sharp rank bounds and all-parameter\\tree--network separation}
\author{Anonymous}
\date{4 October 2026\\\small Candidate manuscript \quad DOI: \ReleaseDOI}
\begin{document}
\maketitle
\begin{abstract}
A generic identifiability theorem can leave exceptional positive parameters at
which a phylogenetic network and a tree induce the same site-pattern
probabilities. We establish all-parameter separation results for two
single-reticulation models. First, for a $d$-state four-cycle with positive
invertible Markov channels, two quartet flattenings have rank $d^2$ and the
third has rank at least $d(d+1)/2$. The latter bound is sharp for every $d\ge2$,
including positive continuous-time reversible examples. Consequently no such
distribution is a mixture of at most $\lfloor d/2\rfloor$ general-Markov
components sharing one tree topology; for nucleotide data this excludes two
components. Bayes rerooting and marginalization extend the result to any
single-reticulation network with at least four observable cycle attachments,
without stationarity or reversibility. Second, a triangle evolving under one
stationary irreducible reversible generator is separated from the entire
closure of the homogeneous stationary reversible tree model for $d\ge3$,
and also for two states with nonuniform stationary frequencies. The proof
identifies the competing generator from an unaffected pair, then combines
strict spectral concavity with a nonuniform equal-input moment obstruction.
Its uniform four-state specialization is a known Jukes--Cantor inequality,
which we explicitly credit. The remaining binary symmetric case admits an
exact tree-mimicry criterion. Finally, the rank bounds yield a computable
pointwise distance certificate and a conservative independent-site test.
These conclusions concern the specified models, not unrestricted network
identifiability or arbitrary rate heterogeneity.
\end{abstract}
\noindent\textbf{Keywords:} phylogenetic networks; general Markov model;
time reversibility; flattening rank; mixture models; identifiability.

\section{Introduction}

A phylogenetic tree describes divergence through a sequence of branching
events. A network can additionally represent reticulation, with different
ancestral routes contributing to the distribution of an observed character.
In a displayed-tree substitution model, one parent is chosen at each
reticulation and evolution follows the resulting tree. Edges retained in
several displayed trees carry the same transition matrices. This sharing is
part of the model, rather than a restriction imposed after fitting an
arbitrary mixture \cite{Gross2021,Sullivant2026}.

The statistical question considered here is whether the complete
site-pattern distribution can distinguish such a network from a tree.
Generic identifiability settles this question outside an exceptional
parameter set of lower dimension. An all-parameter result instead excludes
every coincidence in a specified domain. These are different assertions:
a distribution satisfying a special symmetry can be nongeneric while all
its transition probabilities, branch lengths and inheritance weights remain
strictly positive.

We study two complementary settings. A four-cycle permits a direct
linear-algebraic argument under arbitrary positive invertible general-Markov
(GM) channels. A triangle needs a different argument: we assume one
stationary reversible generator throughout the network and allow a competing
tree to choose its own stationary distribution and generator. The results
are stated separately because the triangle assumption cannot be removed
by the four-cycle proof.

The main rank theorem gives the following sharp universal statement. For a
$d$-state positive Markov four-cycle, the three quartet flattening ranks are
\[
 d^2,\qquad d^2,\qquad r\ge \frac{d(d+1)}2.
\]
The third rank can equal its lower bound, so replacing it by $d^2$ would be
false. The proof uses a determinant factorization for the first two splits
and a filtration of the transformed row space for the third. Since a
mixture of $m$ GM components on one tree topology has rank at most $md$ at
that topology's split, the theorem excludes every such mixture when
$m\le\lfloor d/2\rfloor$. For $d=4$, all three network ranks are at least
10, whereas a two-component single-topology mixture has one rank at most
8. This conclusion permits unrelated parameters in the competing mixture's
components.

The reversible triangle theorem treats every possible spectrum of the
shared generator. A matrix logarithm identifies the competing tree's
generator up to scale. If there are two distinct nonzero decay rates,
strict concavity of a two-route log mixture produces a contradiction.
If all nonzero rates coincide, the generator is an equal-input generator;
a centered third moment then separates the models. This second case is
necessary for an all-parameter assertion. The argument works for every
positive stationary distribution with at least three states. It also
works for two states unless their frequencies are equal. We classify that
last case explicitly and exhibit a positive triangle whose entire tensor
equals that of a positive tree.

\subsection{Relation to previous results}

The use of flattening ranks is established in mathematical phylogenetics.
Casanellas and Fern\'andez-S\'anchez give GM tree-clade rank constraints on
networks \cite{CFS2021}, and develop the corresponding equivariant framework
\cite{CFS2025}. Snyman, Fox and Bryant prove an all-parameter tree rank
formula under positivity and nonsingularity assumptions, with explicit
rerooting and restriction arguments \cite{Snyman2023}. Sullivant gives
universal network rank upper bounds and generic lower bounds obtained from
displayed-tree boundary points \cite[Theorems 7.8 and 7.11]{Sullivant2026}.
The lower-bound argument uses semicontinuity, so it does not itself exclude
positive exceptional parameters. A related five-author theorem, publicly
announced in 2025, likewise states a generic network lower bound
\cite[slide 23]{SullivantIMSI2025}. We compare that public statement; the
complete manuscript of the announced collaboration was unavailable.

All-parameter separation already occurs in more restricted substitution
models. Englander and coauthors obtain a Jukes--Cantor trinet inequality
\cite[Proposition 2.26]{Englander2025}; Currie and coauthors reproduce it
and describe the semialgebraic JC triangle model
\cite[Theorem 3.1 and Appendix A]{Currie2026}. Brits and coauthors prove
full results, modulo triangle orientation, for level-1 networks under JC,
K2P and K3P, together with a broader JC tree--network theorem excluding
2-blobs \cite{Brits2026}. Generic identifiability
of triangle-free level-1 networks under these group-based models provides
an earlier baseline \cite[Theorem 2]{Gross2021}.

Our equal-input calculation specializes exactly to the known JC
obstruction. We do not claim a new JC invariant. The triangle contribution
is the extension to nonuniform equal-input frequencies and its combination
with the exhaustive shared-generator spectral argument. Equal-input models
with arbitrary state spaces and frequencies are themselves well established
\cite{CasanellasSteel2017}. Likewise, the binary classification uses the
standard CFN tripod moment inversion \cite{Hill2024}.

For mixture comparisons, Allman, Rhodes and Sullivant already show that a
limited number of same-topology continuous-time components cannot mimic an
unmixed different tree, subject to its refinement qualification, and obtain
stronger arbitrary-class conclusions
under JC and K2P \cite[Theorems 1 and 3]{ARS2012}. The bound $md$ on a
single-topology mixture flattening is therefore an established tool. The
additional ingredient here is the universal GM four-cycle lower bound.
The JC/K2P specialization of the mixture conclusion is covered more
strongly by the earlier linear-test result.

Table~\ref{tab:comparison} summarizes the closest comparisons. These
comparisons concern hypotheses and quantifiers; they do not assert priority
relative to unavailable work.

\begin{longtable}{@{}p{.23\textwidth}p{.33\textwidth}p{.36\textwidth}@{}}
\caption{Theorem-level comparison. ``Universal'' always refers to the
stated parameter domain.}\label{tab:comparison}\\
\toprule
Result & Model and scope & Relation to the present result\\
\midrule
\endfirsthead
\toprule Result & Model and scope & Relation to the present result\\\midrule
\endhead
Englander et al.; Currie et al.\ \cite{Englander2025,Currie2026}
& Universal JC triangle/trinet obstruction; semialgebraic triangle analysis
& Exact precedent for the uniform four-state moment identity.\\[4pt]
Brits et al.\ \cite{Brits2026}
& Full level-1 results modulo triangle orientation for JC/K2P/K3P; broader JC separation excluding 2-blobs
& Different substitution assumptions; includes more network topologies.\\[4pt]
Casanellas--Fern\'andez-S\'anchez\ \cite{CFS2021,CFS2025}
& Universal GM and equivariant tree-clade rank upper bounds
& Does not supply the positive four-cycle lower bound.\\[4pt]
Sullivant and announced collaboration\ \cite{Sullivant2026,SullivantIMSI2025}
& Universal cut upper bounds and generic network lower bounds
& Does not rule out positive exceptional rank loss.\\[4pt]
Snyman--Fox--Bryant\ \cite{Snyman2023}
& Universal positive nonsingular GM tree rank theorem
& Supplies established tree-rank and rerooting mechanisms.\\[4pt]
Allman--Rhodes--Sullivant\ \cite{ARS2012}
& Universal limited-class mixture versus unmixed-tree results with a refinement qualification; stronger JC/K2P results
& Rank subadditivity is known; the GM network lower bound is the extra input.\\[4pt]
Theorem~\ref{thm:rank}
& Positive invertible GM four-cycle, every $d\ge2$
& Sharp universal crossed rank $d(d+1)/2$ and bounded-mixture exclusion.\\[4pt]
Theorems~\ref{thm:triangle} and \ref{thm:triangle-closure}
& Shared stationary reversible generator; $d\ge3$, or nonuniform binary frequencies
& Separation from the whole competing tree closure, including repeated spectra.\\
\bottomrule
\end{longtable}

\section{Models and elementary rank facts}\label{sec:models}

All state spaces are $[d]=\{1,\ldots,d\}$ with $d\ge2$. A stochastic matrix
is row-stochastic. A \emph{positive} matrix has every entry strictly
positive, and \emph{invertible} always means invertible as a real matrix.
A probability vector is written as a column; its transpose is used for
row-vector propagation. All trees and networks here are finite; each
reticulation has exactly two incoming edges. The GM model on a rooted tree assigns a probability
vector to the root and a stochastic matrix to each directed edge. The joint
law is the root probability multiplied by the edge transition probabilities,
and a leaf tensor is obtained by summing over internal states.

\begin{definition}[Displayed-tree mixture]
At a reticulation with parents $u,v$, choose the incoming edge from $u$ with
probability $\gamma$ and that from $v$ with probability
$\delta=1-\gamma$. Conditional on the choice, the character evolves on the
retained tree. All edges common to the displayed trees retain the same
matrices. Throughout the separation results $0<\gamma<1$.
\end{definition}

The corresponding conditional factor at the reticulation is
$\gamma U_{jl}+\delta V_{kl}$, where $j,k$ are its parental states and $l$
is its own state. Thus the displayed-tree model can also be written as a
DAG model with a constrained conditional distribution
\cite{Sullivant2026}. We do not incorporate a coalescent process.

\begin{definition}[Observable cycle attachment]
For a network with one underlying undirected cycle, delete the cycle edges.
A component attached to a cycle vertex is observable if it contains an
observed leaf. The attachment below the reticulation is required to be
observable. Suppressing unused degree-two cycle vertices gives its reduced
cycle. A triangle has three retained attachment vertices; a four-cycle has
four.
\end{definition}

For a four-leaf tensor $P=(P_{ijkl})$, its flattening
$F_{ij\mid kl}(P)$ is the $d^2\times d^2$ matrix with row index $(i,j)$,
column index $(k,l)$ and entry $P_{ijkl}$, using lexicographic order.
The other two flattenings are defined similarly. Flattening preserves the
Euclidean norm of all tensor entries as the Frobenius matrix norm.

\begin{lemma}[Tree rank and leaf actions]\label{lem:tree-rank}
A $d$-state GM quartet tree has a flattening of rank at most $d$ at its
internal split. A star tree has this bound at all three splits. A mixture
of $m$ GM distributions on one quartet topology has rank at most $md$ at
that topology's split. Applying invertible stochastic channels to the four
leaves preserves all three flattening ranks.
\end{lemma}
\begin{proof}
Cutting an internal tree edge and conditioning on its endpoint states
factors the corresponding flattening as $L^T D_p M R$, where $L,R$ are
conditional leaf-pattern matrices and the middle factors are
$d\times d$. Its rank is at most $d$. A central-state conditioning gives
the star-tree assertion. The mixture statement follows by rank
subadditivity, with no equality of numerical parameters required between
components. These are standard flattening facts; see, for example,
\cite{CFS2021,ARS2012,Snyman2023}.

For leaf channels $E_1,\ldots,E_4$, the first flattening transforms as
\[
 F_{12\mid34}(P')=(E_1^T\otimes E_2^T)
                       F_{12\mid34}(P)(E_3\otimes E_4).
\]
Both multiplying matrices are invertible, as are their analogues for the
other splits. Their inverses need not be stochastic.
\end{proof}

A continuous-time generator $Q$ satisfies $Q\one=0$ and has nonnegative
off-diagonal entries. It is irreducible when its directed rate graph is
strongly connected. If $Q$ is irreducible, $e^{tQ}$ is positive and
invertible for each finite $t>0$. Reversibility with respect to a positive
probability vector $\pi$ means $D Q=Q^T D$, where $D=\diag\pi$.
Then $H=D^{1/2}QD^{-1/2}$ is symmetric. Stationarity means the root
law is $\pi$. A \emph{homogeneous stationary reversible model} uses one
such generator throughout an object, with edge-specific times that are finite
nonnegative real numbers. Cycle-edge times are strictly positive where stated. For $d=4$
this is the usual shared-generator GTR setting. The GM rank results below
require none of these additional generator restrictions.

\section{A sharp universal four-cycle rank bound}\label{sec:rank}

Place the middle ordinary vertex at $i$, the parental vertices at $j,k$,
and the reticulation at $l$. The cycle is $i-j-l-k-i$. Its core tensor is
\begin{equation}\label{eq:core}
 C_{ijkl}=p_iS_{ij}T_{ik}\bigl(\gamma U_{jl}+\delta V_{kl}\bigr),
 \qquad\delta=1-\gamma.
\end{equation}
Conditional on $i$, the two parental states evolve independently through
$S,T$. The reticulation then chooses one of $U,V$. Four pendant channels
map the core states to the observed leaves.

\begin{theorem}[Sharp positive Markov four-cycle ranks]\label{thm:rank}
Let $p$ be strictly positive, let $S,T,U,V$ be positive invertible
$d\times d$ stochastic matrices, and let $0<\gamma<1$. Apply any four
invertible stochastic pendant channels to \eqref{eq:core}, obtaining $P$.
Then
\begin{align}
 \rank F_{ij\mid kl}(P)&=d^2,\qquad
 \rank F_{ik\mid jl}(P)=d^2,\label{eq:adjacent-ranks}\\
 \frac{d(d+1)}2\ &\le\rank F_{il\mid jk}(P)\le d^2.
 \label{eq:crossed-rank}
\end{align}
The lower bound in \eqref{eq:crossed-rank} is attained for every $d\ge2$,
even by positive continuous-time reversible channels.
\end{theorem}

\begin{proof}
Lemma~\ref{lem:tree-rank} allows us to work with $C$.
For the first flattening let $D_0=\diag(p_iS_{ij})$, with diagonal index
$(i,j)$. For $1\le k\le d$ let
$R_k=\one V_{k\cdot}$, a $d\times d$ matrix all of whose rows are row $k$
of $V$. Direct entrywise multiplication gives
\begin{equation}\label{eq:det-factor}
 F_{ij\mid kl}(C)=D_0(T\otimes U)
                 \left[\gamma I_{d^2}+
                       \delta\diag(R_1,\ldots,R_d)\right].
\end{equation}
Indeed the first term has entry
$p_iS_{ij}T_{ik}\gamma U_{jl}$. For the second term,
$U\one=\one$ makes its entry $p_iS_{ij}T_{ik}\delta V_{kl}$.
Since $V_{k\cdot}\one=1$, $R_k$ is a rank-one idempotent.
Each block $\gamma I_d+\delta R_k$ has eigenvalue $\gamma+\delta=1$
once and eigenvalue $\gamma$ with multiplicity $d-1$. Hence
\begin{equation}\label{eq:first-determinant}
 \det F_{ij\mid kl}(C)=
 \left(\prod_{i,j}p_iS_{ij}\right)
 (\det T)^d(\det U)^d\gamma^{d(d-1)}\ne0.
\end{equation}
Exchanging the two parents yields
\begin{equation}\label{eq:second-determinant}
 \det F_{ik\mid jl}(C)=
 \left(\prod_{i,k}p_iT_{ik}\right)
 (\det S)^d(\det V)^d\delta^{d(d-1)}\ne0.
\end{equation}
These formulas use the stated lexicographic index orders.

For the crossed split, view row $(i,l)$ as a matrix in indices $(j,k)$.
Apply to every row matrix the same invertible linear map
\[
 \mathcal L(M)=S^{-T}MT^{-1}.
\]
Let $e_i$ denote coordinate vector $i$ and define
\begin{equation}\label{eq:AiBi}
 A_i=S^{-T}\diag(S_{i1},\ldots,S_{id})U,
 \qquad B_i=T^{-T}\diag(T_{i1},\ldots,T_{id})V.
\end{equation}
The transformed row is
\begin{equation}\label{eq:row-star}
 R_{il}=p_i\{\gamma A_i[:,l]e_i^T+
                   \delta e_iB_i[:,l]^T\}.
\end{equation}
In particular it is supported on row $i$ or column $i$.
Every $A_i$ is invertible: all entries in row $i$ of $S$ are nonzero.
Row-stochasticity gives $A_i\one=B_i\one=e_i$, so
\begin{equation}\label{eq:diagonal-aggregate}
 \sum_l R_{il}=p_i e_ie_i^T.
\end{equation}
Thus the full diagonal subspace $\mathcal D$, of dimension $d$, lies in
the transformed row space.

Put $\mathcal V_i=\operatorname{span}\{R_{il}:l\in[d]\}$ and
\[
 \mathcal W_0=\mathcal D,\qquad
 \mathcal W_i=\mathcal D+\mathcal V_1+\cdots+\mathcal V_i.
\]
For $1\le i\le d-1$, project onto the entries $(h,i)$ with $h>i$.
This projection annihilates $\mathcal W_{i-1}$: diagonal matrices vanish
there, and every earlier family $r<i$ is supported only on row $r$ or
column $r$. Its image on $\mathcal V_i$ has dimension $d-i$.
In fact \eqref{eq:row-star} shows that the image matrix is precisely the
$h>i$ row submatrix of $p_i\gamma A_i$, which has full row rank because
$A_i$ is invertible. Therefore
\[
 \dim\mathcal W_i\ge\dim\mathcal W_{i-1}+d-i.
\]
Summing gives
\[
 \rank F_{il\mid jk}(C)\ge\dim\mathcal W_{d-1}
 \ge d+\sum_{i=1}^{d-1}(d-i)=\frac{d(d+1)}2.
\]
The upper bound is the matrix size.

For sharpness take $T=S$, $V=U$ and $\gamma=\delta=1/2$.
Then \eqref{eq:core} is unchanged by interchanging $j$ and $k$.
The crossed columns $(j,k)$ and $(k,j)$ agree, leaving at most
$d(d+1)/2$ distinct columns. The lower bound forces equality.
For example, with uniform $p$ one can use
\begin{equation}\label{eq:sharp-example}
 S=T=U=V=(1-\epsilon)I+\epsilon\one\one^T/d,
 \qquad 0<\epsilon<1.
\end{equation}
These are positive finite-time transitions of a symmetric irreducible
generator. Positive invertible pendant channels may be added freely.
\end{proof}

\begin{remark}
The crossed rank is generically $d^2$: at an inheritance endpoint the core
reduces to a tree, and its crossed flattening is full rank under the stated
positive channels; continuity of its determinant preserves this rank for
nearby interior weights. This observation agrees with the generic
lower-bound mechanism in \cite{Sullivant2026}. It does not replace the
universal lower bound. We make no assertion that every integer between
$d(d+1)/2$ and $d^2$ occurs.
\end{remark}

\begin{corollary}[Exclusion of bounded single-topology mixtures]
\label{cor:mixture}
A tensor in Theorem~\ref{thm:rank} cannot be a mixture of at most
$\lfloor d/2\rfloor$ $d$-state GM components sharing any one quartet tree
topology. Components may have independent numerical parameters and root
distributions, including degenerate stochastic parameters.
\end{corollary}
\begin{proof}
Such an $m$-component mixture has a split of rank at most $md$ by
Lemma~\ref{lem:tree-rank}, whereas every network split has rank at least
$d(d+1)/2>md$. A star is covered by the same lemma. Nonbinary topologies
can also be resolved by identity channels.
\end{proof}

The common-topology qualification is essential. The four-cycle itself is a
two-component mixture of different displayed topologies. The corollary is
therefore not a claim against arbitrary two-tree mixtures.

\section{Reduction of larger general-Markov networks}\label{sec:gm-large}

We next lift the four-cycle theorem without imposing reversibility. The
underlying rerooting operation is standard \cite[Proposition 6]{Snyman2023};
the issue to check is that it is performed on a common backbone law and
therefore preserves the shared-edge mixture structure.

\begin{lemma}[Four-attachment reduction]\label{lem:gm-reduction}
Let $N$ be a rooted network with exactly one reticulation and one underlying
cycle. Assume a positive root distribution, positive invertible stochastic
edge matrices and interior inheritance probability. Suppose at least four
cycle-attachment components are observable, including the component below
the reticulation. Then a four-leaf marginal of its distribution has the
form of Theorem~\ref{thm:rank}.
\end{lemma}
\begin{proof}
Let $h$ be the reticulation and delete its two incoming edges. The component
containing the original root is a Markov tree, the \emph{backbone}, whose
law is identical under both inheritance choices. The other component is a
descendant tree rooted at $h$. Select one leaf there and one leaf in each
of three distinct ordinary cycle attachments. Deleting $h$ from the cycle
leaves a path in the backbone. Write the three selected attachment vertices
as $j,i,k$ in their path order, with $i$ the middle vertex.

Reroot the backbone at $i$. For completeness, an edge $u\to v$ with
transition $M$ and positive marginal vectors $p_u,p_v$, where
$p_v^T=p_u^TM$, reverses to
\begin{equation}\label{eq:bayes-reverse}
 M^{\mathrm{rev}}=D_{p_v}^{-1}M^TD_{p_u}.
\end{equation}
It is positive, stochastic and invertible, and
$D_{p_u}M=(D_{p_v}M^{\mathrm{rev}})^T$ preserves the edge joint law.
Successive reversals preserve the whole backbone distribution. No incoming
reticulation edge is reversed, and the backbone is rerooted only once,
before multiplying by its unchanged reticulation conditional factor.

Marginalize branches containing no chosen leaf and contract degree-two
paths. Products of positive invertible stochastic matrices retain all
three properties. The paths from $i$ to $j,k$ give $S,T$. Continue from
$j$ along its remaining cycle arm to its parental endpoint and then into
$h$; its product gives $U$. The other arm gives $V$. The selected leaf
paths give four invertible pendant channels. The resulting core law is
exactly \eqref{eq:core}, with a positive marginal distribution at $i$.
Both parental routes remain and the inheritance weight is unchanged.
\end{proof}

\begin{theorem}[Arbitrary-taxon GM separation]\label{thm:gm-network}
A network satisfying Lemma~\ref{lem:gm-reduction} cannot equal a mixture
of at most $\lfloor d/2\rfloor$ $d$-state GM components sharing one tree
topology on its full observed leaf set. In particular it cannot equal any
single GM tree distribution.
\end{theorem}
\begin{proof}
Equality of full distributions implies equality of the four-leaf marginals
selected by the lemma. Restricting each component of a fixed tree topology
produces GM distributions on the same induced quartet topology, although
their numerical parameters may differ. Corollary~\ref{cor:mixture} gives
a contradiction.
\end{proof}

\begin{remark}[Several reticulations]
The preceding proof is for one reticulation. Pruning another cycle can
produce a convex combination of path matrices, which need not be
invertible even if both matrices are positive and invertible. For example
\[
 A=\begin{pmatrix}3/4&1/4\\1/4&3/4\end{pmatrix},\qquad
 B=\begin{pmatrix}1/4&3/4\\3/4&1/4\end{pmatrix}
\]
have this property, whereas $(A+B)/2$ has rank one. This is an obstruction
to that automatic reduction argument, not a claimed counterexample to
every possible multiple-cycle theorem.
\end{remark}

\section{Homogeneous reversible triangles}\label{sec:triangle}

For the remainder of the triangle argument fix a positive probability
vector $\pi$ and an irreducible generator $Q$ reversible with respect to
$\pi$. The same $Q$ acts on every network edge. Let $a,b,c>0$ and
$0<\gamma<1$. The triangle core is
\begin{equation}\label{eq:triangle-core}
 P^0_{ijl}=\pi_i(e^{aQ})_{ij}
       \{\gamma(e^{bQ})_{il}+\delta(e^{cQ})_{jl}\}.
\end{equation}
Attach pendant transitions $e^{xQ},e^{yQ},e^{zQ}$, where $x,y,z\ge0$
are finite. Allowing zero pendant times strengthens the statement; the
ordinary cycle time $a$ and the two parental route times remain positive.
All one-leaf marginals equal $\pi$.

A competing homogeneous stationary reversible tree may choose its own
positive stationary distribution, irreducible generator and finite
nonnegative lengths. By reversibility it can be rooted at the central
vertex of its induced three-leaf tree.

\begin{lemma}[Identification of a competing generator]\label{lem:generator}
If the triangle tensor equals a homogeneous stationary reversible tree
tensor, then the two generators are positive scale multiples of each
other. The scale can be absorbed into the tree lengths.
\end{lemma}
\begin{proof}
One-leaf marginals force the tree stationary distribution to be $\pi$.
Write $D=\diag\pi$ and $L=x+a+y>0$. The unaffected pair has joint matrix
$De^{LQ}$. For a competing generator $Q'$ and tree path length $L'$,
equality implies
\[
 e^{LQ}=e^{L'Q'}.
\]
The left side is positive and is not the identity, so $L'>0$.
Conjugating by $D^{1/2}$ produces exponentials of real symmetric matrices.
Their common value is symmetric positive definite and has a unique real
symmetric logarithm. Hence $LQ=L'Q'$. This argument requires no simple
spectrum and asserts uniqueness only within real symmetric logarithms.
\end{proof}

\begin{theorem}[All-parameter reversible triangle separation]
\label{thm:triangle}
The tensor defined by \eqref{eq:triangle-core} and its pendant channels
cannot equal any homogeneous stationary reversible tree tensor if either
\begin{enumerate}[label=(\roman*)]
\item $d\ge3$; or
\item $d=2$ and $\pi\ne(1/2,1/2)$.
\end{enumerate}
The competing tree may initially choose its own generator and stationary
distribution, and may have finite nonnegative edge lengths.
\end{theorem}
\begin{proof}
By Lemma~\ref{lem:generator}, it suffices to compare a tree using $Q$.
The symmetric conjugate of $Q$ has one zero eigenvalue and $d-1$ strictly
negative eigenvalues. Write the nonzero decay rates as positive numbers.
We treat their exhaustive dichotomy.

\emph{At least two distinct decay rates.}
For a decay rate $t>0$, the first hybrid-involving pair has an attenuation
factor $e^{-(x+z)t}(\gamma e^{-tb}+\delta e^{-t(a+c)})$.
A single tree path would require the same effective extra length
\begin{equation}\label{eq:effective-length}
 F(t)=-\frac1t\log(\gamma e^{-tb}+\delta e^{-t(a+c)})
\end{equation}
at every nonzero decay rate. Put $h(t)=tF(t)$. It has $h(0)=0$, and when
$b\ne a+c$,
\begin{equation}\label{eq:strict-concavity}
 h''(t)=-\frac{\gamma\delta(b-a-c)^2e^{-t(b+a+c)}}
 {\{\gamma e^{-tb}+\delta e^{-t(a+c)}\}^2}<0.
\end{equation}
For $0<r<s$, strict concavity at
$r=(r/s)s+(1-r/s)0$ gives $h(r)/r>h(s)/s$. Thus equality with a tree
forces $b=a+c$. Applying the same argument to the other hybrid-involving
pair forces $a+b=c$. Together these imply $a=0$, a contradiction.
Repeated eigenvalues other than the two chosen distinct values cause no
problem.

\emph{One common nonzero decay rate.}
Suppose all nonzero eigenvalues equal $-\lambda$, with $\lambda>0$.
The reversible spectral decomposition and the one-dimensional stationary
eigenspace give
\begin{equation}\label{eq:equal-input}
 Q=\lambda(\one\pi^T-I).
\end{equation}
This is the established equal-input model, known as F81 for four states
\cite{CasanellasSteel2017}. Every centered state function is an
eigenfunction with attenuation $e^{-\lambda t}$.

Under either hypothesis of the theorem, choose a state $m$ with
$\pi_m\ne1/2$. For $d\ge3$ a state of minimum frequency suffices.
Define
\begin{equation}\label{eq:contrast}
 \phi(i)=\frac{\mathbf1_{\{i=m\}}-\pi_m}
                  {\sqrt{\pi_m(1-\pi_m)}},\qquad
 \kappa=\E_\pi\phi^3=
       \frac{1-2\pi_m}{\sqrt{\pi_m(1-\pi_m)}}\ne0.
\end{equation}
Then $\E_\pi\phi=0$ and $\E_\pi\phi^2=1$. Write
\[
 A=e^{-\lambda a},\quad B=e^{-\lambda b},\quad C=e^{-\lambda c},
 \qquad X=e^{-\lambda x},\quad Y=e^{-\lambda y},\quad Z=e^{-\lambda z}.
\]
For observed leaves $L_1,L_2,L_3$, put
$M_{rs}=\E[\phi(L_r)\phi(L_s)]$ and
$W=\E[\phi(L_1)\phi(L_2)\phi(L_3)]$.
Conditional expectations along an edge multiply a centered function by its
attenuation. This gives
\begin{align}
 M_{12}&=XYA, &
 M_{13}&=XZ(\gamma B+\delta AC), &
 M_{23}&=YZ(\gamma AB+\delta C),\label{eq:triangle-pairs}\\
 W&=XYZ\kappa A(\gamma B+\delta C).&&&\label{eq:triangle-third}
\end{align}
To check the third moment directly, under the first inheritance choice
root the core at $i$. Conditional independence gives
$\E[\phi(i)\phi(j)\phi(l)]=AB\E_\pi\phi^3$.
Under the second choice reroot at $j$ by reversibility, giving $AC\kappa$.
Mixing and multiplying the three pendant factors proves
\eqref{eq:triangle-third}.

On an equal-input tree with pendant attenuation factors $u,v,w$,
conditioning on the central state instead gives
\[
 M_{12}=uv,\qquad M_{13}=uw,\qquad M_{23}=vw,
 \qquad W=\kappa uvw.
\]
Therefore every such tree satisfies
$W^2/\kappa^2=M_{12}M_{13}M_{23}$, including when an edge is zero.
For the triangle, substitution and expansion give
\begin{align}
 \frac{W^2}{\kappa^2}-M_{12}M_{13}M_{23}
 &=(XYZ)^2\bigl[A^2(\gamma B+\delta C)^2
  -A(\gamma B+\delta AC)(\gamma AB+\delta C)\bigr]\notag\\
 &=-(XYZ)^2 A\gamma\delta BC(1-A)^2<0.\label{eq:defect}
\end{align}
Here all displayed factors are positive and $A<1$. This contradiction
completes the proof.
\end{proof}

\subsection{The known JC specialization}

For uniform four-state frequencies, \eqref{eq:defect} is the known JC
triangle obstruction of Englander et al.\ and Currie et al.
\cite{Englander2025,Currie2026}. The equivalence can be stated explicitly.
Relabel Currie's hybrid leaf 1 as our leaf 3, and their leaves 2,3 as our
leaves 1,2. In their absorbed Appendix~A parameters set
\[
 a_2=X,\qquad a_3=Y,\qquad a_5=A,\qquad
 a_4=Z\gamma B,\qquad a_6=Z\delta C.
\]
Their Fourier coordinates $(q_{ACC},q_{CAC},q_{CCA},q_{CGT})$ become
$(M_{12},M_{23},M_{13},W/\kappa)$ in this specialization. Their positive
factorization is exactly the negative of \eqref{eq:defect}. Relabelling
and sign reversal do not yield a new identity. The present use of a
nonuniform contrast and the preceding spectral case have different
hypotheses from that known specialization.

\begin{lemma}[Three-attachment reversible reduction]
\label{lem:triangle-reduction}
A homogeneous stationary reversible single-reticulation network with at
least three observable cycle attachments, including the hybrid attachment,
and positive finite cycle lengths has a triangle marginal of the form
\eqref{eq:triangle-core} with $a,b,c>0$ and unchanged interior inheritance.
\end{lemma}
\begin{proof}
Select one leaf below the reticulation and one in each of two other
observable attachments. Marginalization commutes with the two-term
mixture. Remove unused subtrees and suppress degree-two paths, using
$e^{sQ}e^{tQ}=e^{(s+t)Q}$. If suppression passes through a split vertex,
stationary reversibility gives the fork identity
\begin{equation}\label{eq:fork}
 \sum_i\pi_i(e^{sQ})_{ij}(e^{tQ})_{ik}
       =\pi_j(e^{(s+t)Q})_{jk}.
\end{equation}
Indeed detailed balance changes
$\pi_i(e^{sQ})_{ij}$ to $\pi_j(e^{sQ})_{ji}$, after which matrix
multiplication gives the right side. Distinct retained attachments leave
three nonempty cycle paths, hence positive summed lengths. Both routes
into the hybrid attachment remain. The selected leaf paths give the
pendant lengths, and the weight is unchanged.
\end{proof}

\begin{corollary}\label{cor:reversible-large}
Theorem~\ref{thm:triangle} applies to homogeneous stationary reversible
single-reticulation networks on any number of observed taxa satisfying
Lemma~\ref{lem:triangle-reduction}. For $d=4$ this gives the specified
all-parameter GTR tree--network separation with three or more observable
cycle attachments.
\end{corollary}
\begin{proof}
An equality with a full tree distribution would imply equality of the
selected three-leaf marginals. The induced tree remains homogeneous,
stationary and reversible, so Theorem~\ref{thm:triangle} excludes it.
\end{proof}

\section{The binary symmetric boundary}\label{sec:binary}

The state-space qualification in Theorem~\ref{thm:triangle} is substantive.
When $d=2$ and $\pi=(1/2,1/2)$, every centered contrast has zero third
moment. All irreducible reversible generators with that stationary law
are positive scale multiples of the binary symmetric generator.
The following result substitutes the network's moments into the
established CFN tripod geometry \cite{Hill2024}.

\begin{proposition}[Exact binary symmetric criterion]\label{prop:binary}
In the triangle model with $d=2$ and uniform frequencies, use the
attenuations defined above and put
\[
 q=\gamma B+\delta AC,\qquad r=\gamma AB+\delta C.
\]
The triangle tensor equals a finite positive-length homogeneous stationary
reversible tree tensor if and only if
\begin{equation}\label{eq:binary-criterion}
 Z^2qr<A.
\end{equation}
The unique positive tree attenuation factors are
\begin{equation}\label{eq:binary-inverse}
 u=X\sqrt{Aq/r},\qquad v=Y\sqrt{Ar/q},\qquad
 w=Z\sqrt{qr/A}.
\end{equation}
If $Z^2qr=A$, a finite nonnegative-length tree matches, with its third
edge of length zero. If $Z^2qr>A$, no finite nonnegative-length tree
matches.
\end{proposition}
\begin{proof}
Encode states by spins $\sigma_i\in\{-1,1\}$. Simultaneous sign-reversal
symmetry makes every odd moment zero. Fourier inversion on the three-spin
cube gives
\begin{equation}\label{eq:binary-fourier}
 P(\sigma_1,\sigma_2,\sigma_3)=\frac18
 \{1+\sigma_1\sigma_2M_{12}+\sigma_1\sigma_3M_{13}
                         +\sigma_2\sigma_3M_{23}\}.
\end{equation}
Thus the three pair moments determine the whole tensor. They are
$XYA,XZq,YZr$. A tree must satisfy
$uv=XYA$, $uw=XZq$, $vw=YZr$, whose unique positive solution is
\eqref{eq:binary-inverse}.

The first two factors are strictly below one because
\[
 r-Aq=\delta C(1-A^2)>0,\qquad
 q-Ar=\gamma B(1-A^2)>0,
\]
and $0<X,Y\le1$. All three factors are positive. Positive finite tree
lengths are therefore equivalent to $w<1$, giving
\eqref{eq:binary-criterion}. Allowing zero lengths replaces $<$ by
$\le$. One-leaf marginals exclude a competing nonuniform stationary
law, and the generator scale is absorbed into these factors.
\end{proof}

For an exact positive example choose
\[
 \gamma=\delta=\tfrac12,\qquad A=\tfrac14,
 \qquad B=C=X=Y=Z=\tfrac12.
\]
Then $q=r=5/16$ and the pair moments are
$(1/16,5/64,5/64)$. The matching tree factors are
$(u,v,w)=(1/4,1/4,5/16)$, all strictly below one. At decay rate one the
network lengths are $(\log4,\log2,\log2,\log2,\log2,\log2)$, while
the tree lengths are $(\log4,\log4,\log(16/5))$. The entire common tensor
is
\[
 P(+++) = P(---)=39/256,\qquad
 P(++-) = P(--+)=29/256,
\]
with each remaining pattern equal to $30/256$.
This example concerns a triangle. The GM four-cycle rank theorem applies
even when $d=2$.

\section{Separation from the full reversible tree closure}\label{sec:closure}

Let $\mathcal T_d$ be the set of three-leaf distributions generated by a
homogeneous irreducible reversible generator at stationarity, with a
positive stationary distribution and finite positive tree lengths.
Each competing tree may choose all these parameters independently.
Write $\overline{\mathcal T_d}$ for its Euclidean closure in the
$d^3$-pattern probability simplex. In particular, this closure allows
sequences with changing generators, changing stationary distributions,
unbounded lengths and vanishing edges.

\begin{theorem}[Triangle exclusion from tree closure]
\label{thm:triangle-closure}
Every triangle tensor satisfying either state-space hypothesis of
Theorem~\ref{thm:triangle} lies outside $\overline{\mathcal T_d}$.
\end{theorem}
\begin{proof}
Suppose tree tensors $P_n\in\mathcal T_d$ converge to the triangle tensor
$P$. Their common leaf marginals $\pi_n$ converge to $\pi>0$.
Put $D_n=\diag\pi_n$, $D=\diag\pi$, and let $J_{rs,n},J_{rs}$ be
the pair joint-distribution matrices. Define
\[
 S_{rs,n}=D_n^{-1/2}J_{rs,n}D_n^{-1/2},\qquad
 S_{rs}=D^{-1/2}J_{rs}D^{-1/2}.
\]
These converge pairwise. With $H=D^{1/2}QD^{-1/2}$, the triangle's
normalized matrices are
\begin{align}
 S_{12}&=e^{(x+a+y)H},\notag\\
 S_{13}&=e^{(x+z)H}\{\gamma e^{bH}+\delta e^{(a+c)H}\},\label{eq:pair-spd}\\
 S_{23}&=e^{(y+z)H}\{\gamma e^{(a+b)H}+\delta e^{cH}\}.\notag
\end{align}
Each is symmetric positive definite. Its eigenvalue on $\sqrt\pi$ is
one, and every other eigenvalue lies strictly between zero and one.
Thus the symmetric logarithms
\[
 A_{rs}=-\log S_{rs},\qquad t_{rs}=\tr A_{rs}>0
\]
are finite. Continuity of this logarithm on the positive definite cone
gives $A_{rs,n}\to A_{rs}$ and $t_{rs,n}\to t_{rs}$.

Write the symmetric generator of tree $n$ as
$H_n=D_n^{1/2}Q_nD_n^{-1/2}$ and its three central pendant lengths as
$\ell_{1,n},\ell_{2,n},\ell_{3,n}>0$. Then
\[
 A_{rs,n}=-(\ell_{r,n}+\ell_{s,n})H_n.
\]
Normalize through the unaffected pair:
\begin{equation}\label{eq:normalized-generator}
 B_n=\frac{A_{12,n}}{t_{12,n}}
 =\frac{-H_n}{-\tr H_n}
 \longrightarrow B=\frac{A_{12}}{t_{12}}
 =\frac{-H}{-\tr H}.
\end{equation}
This step absorbs arbitrary generator and length scales; no bound on
$Q_n$ is assumed. Define
\[
 s_{1,n}=\tfrac12(t_{12,n}+t_{13,n}-t_{23,n})
\]
and cyclically define $s_{2,n},s_{3,n}$. They equal
$(-\tr H_n)\ell_{r,n}>0$ and converge to finite values $s_r\ge0$.
The tree edge channels consequently converge to
\begin{equation}\label{eq:limit-channels}
 E_r=D^{-1/2}e^{-s_r B}D^{1/2}
     =e^{\{s_r/(-\tr Q)\}Q}.
\end{equation}
The root distribution converges to $\pi$. The finite sum over the central
state is continuous, so the limiting tree tensor is exactly the tree with
generator $Q$, root law $\pi$ and the finite nonnegative lengths in
\eqref{eq:limit-channels}. This contradicts
Theorem~\ref{thm:triangle}.
\end{proof}

The same closure exclusion holds for a larger single-reticulation network
with a triangle marginal satisfying the theorem. A convergent sequence of
full tree tensors would induce a convergent sequence of their three-leaf
tree marginals, contradicting Theorem~\ref{thm:triangle-closure}.

\begin{corollary}[Uniform margins on compact network domains]
\label{cor:compact}
Let $K$ be any nonempty compact set of triangle parameter tuples satisfying
Theorem~\ref{thm:triangle}. In the binary case, require every tuple to have
nonuniform stationary frequencies. Then
\[
 \epsilon_K=\min_{\theta\in K}\min_{T\in\overline{\mathcal T_d}}
                \TV(P_\theta,T)>0,
 \qquad \TV(p,q)=\tfrac12\sum_\omega|p_\omega-q_\omega|.
\]
The competing tree class remains unrestricted.
\end{corollary}
\begin{proof}
The image of $K$ is compact by continuity of the model. The closed tree
model is a closed subset of a finite probability simplex and is compact.
The distance minimum is attained; if zero, it contradicts
Theorem~\ref{thm:triangle-closure}.
\end{proof}

This corollary gives existence, not an evaluated useful lower bound on a
biological parameter box. The unrestricted open network class has distance
infimum zero from trees: as $\gamma$ approaches an endpoint, the network
approaches one of its displayed trees. Pointwise exclusion from tree
closure is therefore consistent with intersection of the two model
closures.

\section{Observable rank margins and independent-site testing}
\label{sec:statistics}

The rank theorem supplies a directly computable pointwise certificate.
For $1\le m\le\lfloor d/2\rfloor$, let $\mathcal M_m$ be the union,
over quartet tree topologies, of their GM mixtures with at most $m$
components on a common topology. If
$\sigma_1\ge\cdots\ge\sigma_{d^2}$ are a flattening's singular values,
put
\begin{equation}\label{eq:rank-residual}
 R_m(P)=\min_{s\in\{12|34,13|24,14|23\}}
       \left(\sum_{a>md}\sigma_a(F_s(P))^2\right)^{1/2}.
\end{equation}
Tail-singular-value scores and their interpretation as lower bounds on
phylogenetic model distance are established tools
\cite{EckartYoung1936,LongKubatko2021}. The role of
Theorem~\ref{thm:rank} is to make this score positive throughout the
specified network domain.

\begin{corollary}[Computable pointwise margin]\label{cor:residual}
Every $P$ in Theorem~\ref{thm:rank} satisfies
\[
 \dist_2(P,\overline{\mathcal M_m})\ge R_m(P)>0.
\]
The function $R_m$ is $1$-Lipschitz in the Euclidean norm of tensor entries.
\end{corollary}
\begin{proof}
For a split $s$, let $\mathcal D_s$ be the closed determinantal set of
tensors whose $s$-flattening has rank at most $md$, and let
$\mathcal D=\bigcup_s\mathcal D_s$. By Lemma~\ref{lem:tree-rank},
$\mathcal M_m\subseteq\mathcal D$. The finite union $\mathcal D$ is
closed, so it also contains $\overline{\mathcal M_m}$.
Flattening is an isometry and the Eckart--Young theorem gives
\[
 \dist_2(P,\mathcal D)=R_m(P).
\]
All three network ranks exceed $md$, hence $P\notin\mathcal D$ and this
distance is positive. Distance to any nonempty set is $1$-Lipschitz.
\end{proof}

\begin{corollary}[A conservative iid test]\label{cor:test}
Observe $n$ independent identically distributed quartet patterns and let
$\widehat P_n$ be their empirical distribution. Put $K_4=d^4$ and, for
$0<\alpha<1$, define
\[
 \tau_{n,\alpha}=\sqrt{\frac{K_4\log(2K_4/\alpha)}{2n}}.
\]
Reject $P\in\overline{\mathcal M_m}$ when
$R_m(\widehat P_n)>\tau_{n,\alpha}$. Its type-I error is at most
$\alpha$ uniformly over the null. Its power tends to one at every fixed
positive four-cycle distribution satisfying Theorem~\ref{thm:rank}.
\end{corollary}
\begin{proof}
If $\|\widehat P_n-P\|_2>\eta$, some coordinate differs by more than
$\eta/\sqrt{K_4}$. Each coordinate is an empirical mean of independent
Bernoulli indicators. Two-sided Hoeffding \cite{Hoeffding1963} and a union
bound give
\begin{equation}\label{eq:coordinate-hoeffding}
 \Prob_P\{\|\widehat P_n-P\|_2>\eta\}
       \le2K_4e^{-2n\eta^2/K_4}.
\end{equation}
No independence between coordinates is needed. Under the null,
$R_m(P)=0$, so the Lipschitz property bounds $R_m(\widehat P_n)$ by the
norm error. Substitution of $\tau_{n,\alpha}$ proves the level bound.
At each allowed fixed alternative, $R_m(P)>0$, the empirical tensor
converges almost surely to $P$, and the threshold tends to zero.
\end{proof}

For comparison, Corollary~\ref{cor:compact} gives a related existence
statement for triangle testing. Let $K$ and $\epsilon_K$ be as there.
For $n$ iid three-leaf patterns, let $\widehat P_n$ denote their empirical
distribution. Consider the idealized rule selecting the network when
\[
 \dist_{\TV}(\widehat P_n,\overline{\mathcal T_d})>\epsilon_K/2.
\]
Both error probabilities are at most
\begin{equation}\label{eq:triangle-test}
 \min\{1,\ 2^{d^3}e^{-n\epsilon_K^2/2}\}.
\end{equation}
To see this, total variation on an alphabet of size $M=d^3$ equals
$\max_{A\subseteq[M]}(\widehat P_n(A)-P(A))$. Apply one-sided
Hoeffding to each subset and take a union bound over $2^M$ subsets.
If the variation error is less than $\epsilon_K/2$, the rule is correct
under either class by the triangle inequality.

These are conservative theoretical consequences. The triangle rule assumes
a known margin and a distance to the full tree closure that we do not
provide an efficient algorithm to compute. The four-cycle score is
computable by three singular-value decompositions, but a finite-precision
implementation needs numerical error control. Neither rule addresses site
dependence, optimal sample complexity or empirical biological calibration.
For larger networks, any selected quartet supplied by
Lemma~\ref{lem:gm-reduction} can be used with the quartet rule; a
multiple-quartet selection scheme would require its own error analysis.

\section{Examples, limitations and further questions}\label{sec:discussion}

The sharp family \eqref{eq:sharp-example} has crossed ranks $3,6,10$ for
$d=2,3,4$, respectively. It shows that exceptional rank loss can persist
strictly inside the positive continuous-time parameter domain, while the
rank remains above every tree threshold. The determinant formulas and the
row-space filtration explain why this loss is harmless for the stated
separation result. Numerical evidence of generic full rank would not
establish this fact.

The binary example in Section~\ref{sec:binary} gives the opposite type of
boundary information: removing the nonuniform-frequency or $d\ge3$
condition from triangle separation would make that statement false.
The full distribution, not just selected pair summaries, agrees with a
tree there. Conversely, the repeated-spectrum equal-input case is included
in the positive separation theorem whenever a nonzero third centered
moment is available.

The two model classes have distinct strengths. The four-cycle result
permits nonstationarity, arbitrary independent edge matrices and a limited
number of site-mixture classes on the competing tree. It still requires
positive invertible network channels and four observable attachments.
The triangle theorem permits three attachments but requires a shared
reversible generator at stationarity. It does not cover unrelated
edge-specific generators, and its closure theorem concerns the homogeneous
reversible tree class rather than every GM tree. The GM comparator permits
singular and boundary stochastic parameters; the network side does not.

Neither result identifies reticulation orientation, all edge parameters or
an arbitrary network topology. An arbitrary independently parameterized
mixture of different tree topologies is outside the exclusion claim.
Unrestricted rate mixtures and coalescent processes are also outside the
models. The finite-mixture result concerns an upper bound on the number of
components sharing a topology, rather than every possible representation
of site heterogeneity.

Exact separation does not give a uniform practical detection margin on the
whole interior. In \eqref{eq:defect}, the gap tends to zero at inheritance
endpoints, as the ordinary cycle length vanishes, or as pendant attenuation
becomes small. In particular
$(1-e^{-\lambda a})^2\sim\lambda^2a^2$ as $a\to0$.
The same distinction holds for the rank residual: it is positive at every
allowed point but can be arbitrarily small across the model. Compact
parameter restrictions give uniform margins only because the resulting
network image stays away from the competing closed set.

Several extensions remain separate questions. A general description of the
crossed-rank strata could refine the sharp minimum proved here.
Multiple-reticulation reductions require control of effective observation
channels, because convex mixtures can lose invertibility. Triangle
separation under independently varying reversible edge generators is not
settled by the spectral proof. These questions require additional
arguments; the present proofs do not answer them by genericity.

\section*{Declarations}

\paragraph{Data and code availability.}
The results are analytic. The accompanying candidate archive supplies source
files and finite exact verification examples. Such examples check algebra
and implementations; they do not replace the universal proofs or constitute
proof-assistant verification. No empirical biological dataset is used.

\paragraph{Authorship and AI assistance.}
The manuscript is released under the author name Anonymous. Conceptualization,
formal analysis, writing and verification were developed through an
AI-assisted workflow. The manuscript is an unrefereed candidate; internal
AI-assisted checking is not external peer review.

\paragraph{Funding, interests and ethics.}
No funding information or competing interests were declared for this
candidate. No human participants, identifiable personal data or animal
experiments were involved; ethics approval is not applicable.

\paragraph{Reuse.}
The original manuscript, claims and accompanying public records are released
under CC0 1.0 (CC0-1.0). Original executable verification code is released
under the MIT License. Cited works retain their own rights.

\begin{thebibliography}{99}

\bibitem{ARS2012}
E.~S. Allman, J.~A. Rhodes, and S.~Sullivant.
When do phylogenetic mixture models mimic other phylogenetic models?
\emph{Systematic Biology} \textbf{61}(6), 1049--1059 (2012).
\doi{10.1093/sysbio/sys064}.
Inspected preprint: \href{https://arxiv.org/abs/1202.2396v2}{arXiv:1202.2396v2}.

\bibitem{Brits2026}
J.~Brits, N.~Holtgrefe, L.~van Iersel, and S.~Martin.
On tree--network distinguishability and full identifiability of
phylogenetic networks (2026).
\href{https://arxiv.org/abs/2607.12919v3}{arXiv:2607.12919v3}, 25 August 2026.
\doi{10.48550/arXiv.2607.12919}.

\bibitem{CFS2021}
M.~Casanellas and J.~Fern\'andez-S\'anchez.
Rank conditions on phylogenetic networks.
In \emph{Extended Abstracts GEOMVAP 2019}, Trends in Mathematics,
pp.~65--69 (2021). \doi{10.1007/978-3-030-84800-2_11}.
Inspected preprint: \href{https://arxiv.org/abs/2004.12988v1}{arXiv:2004.12988v1}.

\bibitem{CFS2025}
M.~Casanellas and J.~Fern\'andez-S\'anchez.
Phylogenetic invariants: from the general Markov to equivariant models.
\emph{SIAM Journal on Applied Algebra and Geometry} \textbf{9}(4),
765--793 (2025). \doi{10.1137/24M1644821}.
Theorem numbers refer to the inspected accepted manuscript.

\bibitem{CasanellasSteel2017}
M.~Casanellas and M.~Steel.
Phylogenetic mixtures and linear invariants for equal input models.
\emph{Journal of Mathematical Biology} \textbf{74}(5), 1107--1138 (2017).
\doi{10.1007/s00285-016-1055-8}.

\bibitem{Currie2026}
B.~Currie, A.~K. Englander, J.~A. Esparza-Lozano, E.~Gross, M.~Hill,
C.~Long, D.~Olds, K.~O'Connor, U.~Ranasinghe, and C.~Sum.
Semialgebraic conditions for identifying triangles in phylogenetic networks
(2026). \href{https://arxiv.org/abs/2606.26673v1}{arXiv:2606.26673v1},
25 June 2026. \doi{10.48550/arXiv.2606.26673}.

\bibitem{EckartYoung1936}
C.~Eckart and G.~Young.
The approximation of one matrix by another of lower rank.
\emph{Psychometrika} \textbf{1}(3), 211--218 (1936).
\doi{10.1007/BF02288367}.

\bibitem{Englander2025}
A.~K. Englander, M.~Frohn, E.~Gross, N.~Holtgrefe, L.~van Iersel,
M.~Jones, and S.~Sullivant.
Identifiability of phylogenetic level-2 networks under the Jukes--Cantor model.
bioRxiv preprint, version posted 24 April 2025.
\doi{10.1101/2025.04.18.649493}.
Proposition 2.26 in the inspected version.

\bibitem{Gross2021}
E.~Gross, L.~van Iersel, R.~Janssen, M.~Jones, C.~Long, and Y.~Murakami.
Distinguishing level-1 phylogenetic networks on the basis of data generated
by Markov processes.
\emph{Journal of Mathematical Biology} \textbf{83}, Article 32 (2021).
\doi{10.1007/s00285-021-01653-8}.

\bibitem{Hill2024}
M.~Hill, S.~Roch, and J.~I. Rodriguez.
Maximum likelihood estimation for unrooted 3-leaf trees: an analytic
solution for the CFN model.
\emph{Bulletin of Mathematical Biology} \textbf{86}, Article 106 (2024).
\doi{10.1007/s11538-024-01340-x}.

\bibitem{Hoeffding1963}
W.~Hoeffding.
Probability inequalities for sums of bounded random variables.
\emph{Journal of the American Statistical Association} \textbf{58}(301),
13--30 (1963). \doi{10.1080/01621459.1963.10500830}.

\bibitem{LongKubatko2021}
C.~Long and L.~Kubatko.
Hypothesis testing with rank conditions in phylogenetics.
\emph{Frontiers in Genetics} \textbf{12}, 664357 (2021).
\doi{10.3389/fgene.2021.664357}.

\bibitem{Snyman2023}
J.~Snyman, C.~Fox, and D.~Bryant.
Parsimony and the rank of a flattening matrix.
\emph{Journal of Mathematical Biology} \textbf{86}, Article 44 (2023).
\doi{10.1007/s00285-023-01875-y}.

\bibitem{Sullivant2026}
S.~Sullivant.
Phylogenetic network models as graphical models.
\emph{Bulletin of Mathematical Biology} \textbf{88}, Article 136 (2026).
\doi{10.1007/s11538-026-01695-3}.
Theorem numbers refer to
\href{https://arxiv.org/abs/2507.23056v2}{arXiv:2507.23056v2}.

\bibitem{SullivantIMSI2025}
S.~Sullivant.
Phylogenetic network models and graphical models.
IMSI lecture slides, 24 July 2025, slide 23 of 25.
\url{https://cdn.imsi.institute/videos/59496/IKcB160s62/slides.pdf}.
The cited slide attributes its generic rank theorem to M.~Casanellas,
J.~Fern\'andez-S\'anchez, E.~Gross, B.~Hollering and S.~Sullivant.

\end{thebibliography}
\end{document}
